%
\begin{isabellebody}%
\setisabellecontext{Structural{\isacharunderscore}{\kern0pt}Induction}%
%
\begin{isamarkuptext}%
\section*{Proofs by Structural Induction}%
\end{isamarkuptext}\isamarkuptrue%
%
\isadelimtheory
%
\endisadelimtheory
%
\isatagtheory
\isacommand{theory}\isamarkupfalse%
\ Structural{\isacharunderscore}{\kern0pt}Induction\isanewline
\ \ \isakeyword{imports}\ Main\isanewline
\isakeyword{begin}%
\endisatagtheory
{\isafoldtheory}%
%
\isadelimtheory
%
\endisadelimtheory
%
\begin{isamarkuptext}%
Defining a well-formed (recursive) algebraic data type automatically, in Isabelle/HOL, leads
      to the generation of an induction principle. For instance, for the following binary trees data type%
\end{isamarkuptext}\isamarkuptrue%
\isacommand{datatype}\isamarkupfalse%
\ {\isacharprime}{\kern0pt}a\ tree\ {\isacharequal}{\kern0pt}\ Leaf\ {\isacharbar}{\kern0pt}\ Node\ {\isachardoublequoteopen}{\isacharprime}{\kern0pt}a\ tree{\isachardoublequoteclose}\ {\isacharprime}{\kern0pt}a\ {\isachardoublequoteopen}{\isacharprime}{\kern0pt}a\ tree{\isachardoublequoteclose}%
\begin{isamarkuptext}%
that induction principle is 

  \[\frac{P\; \text{Leaf};\;\;\; [\forall\; t_1\; t_2\; a.\; P\; t_1;\; P\; t_2\; \Longrightarrow\; P\; \text{Node t1 a t2}]}{P\; t}.\]

As a general rule of thumb, you should follow the same heuristics for induction: manipulate the goal
until you can apply the Induction Hypothesis (I.H.). In this example, you have two induction hypotheses, one
for each subtree.%
\end{isamarkuptext}\isamarkuptrue%
%
\begin{isamarkuptext}%
The following is an example proof by structural induction on trees.
      After replicating it on your own, take a close look at the two goals: the base case and the step case, and note the two I.H.'s,
      each corresponding to a subtree.

      \emph{Exercise}: write down a detailed manual Isabelle/HOL proof using only forward reasoning
                (i.e.\ \isatt{O{\kern0pt}F{\kern0pt}}), backward reasoning (i.e.\ \isatt{r{\kern0pt}u{\kern0pt}l{\kern0pt}e{\kern0pt}}), and
                substitution (i.e.\ \isatt{s{\kern0pt}u{\kern0pt}b{\kern0pt}s{\kern0pt}t{\kern0pt}}).%
\end{isamarkuptext}\isamarkuptrue%
\isacommand{lemma}\isamarkupfalse%
\ {\isachardoublequoteopen}set\ {\isacharparenleft}{\kern0pt}inorder\ t{\isacharparenright}{\kern0pt}\ {\isacharequal}{\kern0pt}\ tree{\isacharunderscore}{\kern0pt}set\ t{\isachardoublequoteclose}\isanewline
%
\isadelimproof
%
\endisadelimproof
%
\isatagproof
\isacommand{proof}\isamarkupfalse%
{\isacharparenleft}{\kern0pt}induction\ t{\isacharparenright}{\kern0pt}\isanewline
\ \ \isacommand{case}\isamarkupfalse%
\ Leaf\isanewline
\ \ \isacommand{then}\isamarkupfalse%
\ \isacommand{show}\isamarkupfalse%
\ {\isacharquery}{\kern0pt}case\isanewline
\ \ \ \ \isacommand{by}\isamarkupfalse%
\ auto\isanewline
\isacommand{next}\isamarkupfalse%
\isanewline
\ \ \isacommand{case}\isamarkupfalse%
\ {\isacharparenleft}{\kern0pt}Node\ t{\isadigit{1}}\ x{\isadigit{2}}\ t{\isadigit{2}}{\isacharparenright}{\kern0pt}\isanewline
\ \ \isacommand{then}\isamarkupfalse%
\ \isacommand{show}\isamarkupfalse%
\ {\isacharquery}{\kern0pt}case\isanewline
\ \ \ \ \isacommand{by}\isamarkupfalse%
\ auto\isanewline
\isacommand{qed}\isamarkupfalse%
%
\endisatagproof
{\isafoldproof}%
%
\isadelimproof
\isanewline
%
\endisadelimproof
\isanewline
\isanewline
\isanewline
%
\isadelimtheory
\isanewline
%
\endisadelimtheory
%
\isatagtheory
\isacommand{end}\isamarkupfalse%
%
\endisatagtheory
{\isafoldtheory}%
%
\isadelimtheory
%
\endisadelimtheory
%
\end{isabellebody}%
\endinput
%:%file=Structural_Induction.tex%:%
%:%6=1%:%
%:%14=3%:%
%:%15=3%:%
%:%16=4%:%
%:%17=5%:%
%:%26=7%:%
%:%27=8%:%
%:%29=11%:%
%:%30=11%:%
%:%32=15%:%
%:%33=16%:%
%:%34=17%:%
%:%35=18%:%
%:%36=19%:%
%:%37=20%:%
%:%38=21%:%
%:%42=58%:%
%:%43=59%:%
%:%44=60%:%
%:%45=61%:%
%:%46=62%:%
%:%47=63%:%
%:%48=64%:%
%:%50=66%:%
%:%51=66%:%
%:%58=67%:%
%:%59=67%:%
%:%60=68%:%
%:%61=68%:%
%:%62=69%:%
%:%63=69%:%
%:%64=69%:%
%:%65=70%:%
%:%66=70%:%
%:%67=71%:%
%:%68=71%:%
%:%69=72%:%
%:%70=72%:%
%:%71=73%:%
%:%72=73%:%
%:%73=73%:%
%:%74=74%:%
%:%75=74%:%
%:%76=75%:%
%:%82=75%:%
%:%85=76%:%
%:%86=77%:%
%:%87=78%:%
%:%90=79%:%
%:%95=80%:%
